\cleardoublepage
\phantomsection
\begin{planbox}[แผนบริหารการสอนประจำบทที่ 2\quad เวกเตอร์และการเคลื่อนที่ในของเหลว (สัปดาห์ที่ 2)]
\noindent\textbf{1. หัวข้อเนื้อหาประจำบท (บรรยาย 2 ชั่วโมง และปฏิบัติการ 3 ชั่วโมง)}
\begin{itemize}
    \item 2.1 ปริมาณสเกลาร์และปริมาณเวกเตอร์ในระบบชีวภาพ (ระยะทาง การกระจัด อัตราเร็ว และความเร็ว)
    \item 2.2 กลไกการเคลื่อนที่ของแบคทีเรียแบบสลับวิ่งและหมุนสุ่ม (Run and Tumble) และการเคลื่อนที่เข้าหาสารอาหาร
    \item 2.3 การเคลื่อนที่แบบบราวน์และการแพร่ของโมเลกุลในอาหารเลี้ยงเชื้อ
    \item 2.4 ปฏิบัติการที่ 2 การวิเคราะห์กราฟการเคลื่อนที่และการสังเกตการเคลื่อนที่แบบบราวน์เทียบกับการเคลื่อนที่จริงของจุลินทรีย์
\end{itemize}

\noindent\textbf{2. วัตถุประสงค์เชิงพฤติกรรมและผลลัพธ์การเรียนรู้รายวิชา (CLOs ที่เชื่อมโยง PLOs)}
\begin{enumerate}
    \item จำแนกความแตกต่างระหว่างระยะทางกับการกระจัด และคำนวณอัตราเร็วเฉลี่ยของการเคลื่อนที่ของเซลล์ได้ \textbf{(CLO1 เชื่อมโยง PLO2)}
    \item อธิบายเหตุผลทางฟิสิกส์ว่าทำไมการสั่นแบบบราวน์ (Brownian Motion) จึงมีระยะกระจัดลัพธ์เป็นศูนย์ และแยกแยะออกจากการว่ายเคลื่อนที่จริงของแบคทีเรียที่มีแฟลกเจลลาได้ \textbf{(CLO1, CLO2 เชื่อมโยง PLO2, PLO3)}
    \item ปฏิบัติการทดลองวัดอัตราเร็วการเคลื่อนที่และเขียนกราฟระยะทางกับเวลาได้อย่างถูกต้องและปลอดภัย \textbf{(CLO3, CLO4 เชื่อมโยง PLO1, PLO5, PLO6)}
\end{enumerate}

\noindent\textbf{3. กิจกรรมการจัดการเรียนรู้เชิงรุก (Active Learning \& OBE)}
\begin{itemize}
    \item \textbf{Problem-Based Learning (PBL)}\quad ตั้งคำถามกระตุ้นคิด ``เมื่อส่องกล้องจุลทรรศน์แบบหยดแขวน (Hanging Drop) แล้วเห็นเซลล์สั่นดุ๊กดิ๊กอยู่กับที่ แปลว่าแบคทีเรียตัวนั้นว่ายน้ำได้จริงหรือไม่''
    \item \textbf{ฝึกคำนวณด้วยสามเหลี่ยมสูตร $s = v \times t$}\quad คำนวณเวลาที่แบคทีเรียใช้ว่ายข้ามขอบเขตการมองเห็นของกล้องจุลทรรศน์
    \item \textbf{ปฏิบัติการทดลอง}\quad ทดลองจับเวลาการเคลื่อนที่แนวตรง วาดกราฟความชัน และสังเกตการแพร่ของหยดสีในน้ำอุณหภูมิต่างกัน
\end{itemize}

\noindent\textbf{4. สื่อการเรียนรู้และแหล่งเรียนรู้}
\begin{itemize}
    \item เอกสารคำสอนบทที่ 2 แผนภาพเวกเตอร์การกระจัดของแบคทีเรีย และชุดทดลองการเคลื่อนที่และจานเพาะเชื้อผสมวุ้น
\end{itemize}

\noindent\textbf{5. การวัดและประเมินผลประจำบท}
\begin{itemize}
    \item แบบทดสอบย่อยท้ายคาบเรื่องสเกลาร์ เวกเตอร์ และการคำนวณอัตราเร็ว \textbf{(CLO1)}
    \item การตรวจกราฟและรายงานปฏิบัติการที่ 2 \textbf{(CLO3, CLO4)}
\end{itemize}
\end{planbox}
\clearpage

\chapter{เวกเตอร์และการเคลื่อนที่ในของเหลว}
\index{เวกเตอร์}\index{สเกลาร์}\index{การกระจัด}\index{การเคลื่อนที่แบบบราวน์}

\section{ปริมาณสเกลาร์และปริมาณเวกเตอร์ในระบบชีวภาพ}

ในการศึกษาการเคลื่อนที่ของจุลินทรีย์หรือการไหลของของเหลวในห้องปฏิบัติการ ปริมาณทางฟิสิกส์แบ่งออกเป็น 2 ประเภทหลัก ได้แก่ \textbf{ปริมาณสเกลาร์} (\textbf{Scalar Quantity}) คือปริมาณที่บอกเฉพาะขนาดเพียงอย่างเดียวก็เข้าใจสมบูรณ์ เช่น มวล เวลา อุณหภูมิ พื้นที่ผิวโคโลนี และ\textbf{ระยะทาง} (\textbf{Distance}, $s$) ส่วนอีกประเภทหนึ่งคือ \textbf{ปริมาณเวกเตอร์} (\textbf{Vector Quantity}) คือปริมาณที่ต้องระบุทั้งขนาดและทิศทางควบคู่กันเสมอ เช่น แรง สนามไฟฟ้า ความเร็ว และ\textbf{การกระจัด} (\textbf{Displacement}, $\vect{d}$)

ความแตกต่างระหว่างระยะทางกับการกระจัดมีความสำคัญอย่างยิ่งในการสังเกตพฤติกรรมของจุลินทรีย์ภายใต้กล้องจุลทรรศน์ หากแบคทีเรียตัวหนึ่งเคลื่อนที่วกไปวนมาเป็นเส้นทางคดเคี้ยวรวมระยะทางทั้งหมด $s = 120\text{ }\mu\text{m}$ แต่จุดสุดท้ายอยู่ห่างจากจุดเริ่มต้นในแนวเส้นตรงเพียง $d = 30\text{ }\mu\text{m}$ ในเวลา $t = 10\text{ s}$ เราจะคำนวณ\textbf{อัตราเร็วเฉลี่ย} (\textbf{Average Speed}, $v$) และ\textbf{ขนาดความเร็วเฉลี่ย} (\textbf{Average Velocity}, $|\vect{v}|$) ได้จากสมการที่ (2.1)

\begin{equation}
    v = \frac{s}{t} \qquad \text{และ} \qquad |\vect{v}_{\text{avg}}| = \frac{d}{t}
\end{equation}

\section{กลไกการเคลื่อนที่ของแบคทีเรียและการเคลื่อนที่แบบบราวน์}

แบคทีเรียที่มีแฟลกเจลลา (\textbf{Flagella}) เช่น \textit{Escherichia coli} หรือ \textit{Pseudomonas aeruginosa} เคลื่อนที่ในของเหลวด้วยกลไกสลับระหว่าง\textbf{ช่วงพุ่งตรง} (\textbf{Run} เมื่อแฟลกเจลลาหมุนทวนเข็มนาฬิกาเป็นมัดเดียวกัน) กับ\textbf{ช่วงหมุนสุ่มเปลี่ยนทิศ} (\textbf{Tumble} เมื่อแฟลกเจลลาหมุนตามเข็มนาฬิกาแล้วกระจายออก) เมื่อมีสารอาหารดึงดูด (\textbf{Chemotaxis}) แบคทีเรียจะยืดเวลาช่วงพุ่งตรงเข้าหาสารอาหารให้นานขึ้น ทำให้เกิดเวกเตอร์การกระจัดลัพธ์มุ่งไปสู่แหล่งอาหารอย่างชัดเจน ดังแสดงใน\mbox{ภาพที่} 2.1

\begin{figure}[H]
\centering
\resizebox{0.88\textwidth}{!}{%
\begin{tikzpicture}[font=\sffamily, >=Stealth]
    % Left Box: Brownian Motion
    \draw[fill=boxBlueBg, draw=physRoyalBlue, line width=1.2pt, rounded corners=4pt] (0,0) rectangle (6.8,4.6);
    \node[font=\bfseries\color{physRoyalBlue}] at (3.4,4.15) {(ก) การสั่นแบบบราวน์ (Brownian Motion)};
    \draw[physSlate, line width=1.0pt, ->] (3.2,2.1) -- (3.7,2.5) -- (3.1,2.8) -- (2.7,2.2) -- (3.5,1.7) -- (3.0,1.9) -- (3.3,2.2);
    \fill[crimsonAlert] (3.2,2.1) circle (3pt);
    \node[font=\footnotesize, align=center, color=charcoal] at (3.4,0.65) {ถูกโมเลกุลน้ำชนทุกทิศทาง\\ระยะกระจัดลัพธ์ $\vect{d} \approx 0$ (ไม่ได้ว่ายเอง)};

    % Right Box: True Motility (Chemotaxis)
    \draw[fill=boxGreenBg, draw=bioEmerald, line width=1.2pt, rounded corners=4pt] (7.6,0) rectangle (14.8,4.6);
    \node[font=\bfseries\color{bioEmerald!90!black}] at (11.2,4.15) {(ข) การเคลื่อนที่จริงเข้าหาสารอาหาร (Chemotaxis)};
    \fill[physAmber!40] (13.6,3.1) circle (0.75cm);
    \node[font=\footnotesize\bfseries, color=physAmber!90!black] at (13.6,3.1) {สารอาหาร};
    \draw[bioEmerald!90!black, line width=1.2pt, ->] (8.4,1.3) -- (9.8,2.1) -- (9.6,1.6) -- (11.5,2.6) -- (11.3,2.2) -- (13.1,3.0);
    \draw[crimsonAlert, line width=1.6pt, dashed, ->] (8.4,1.3) -- node[below right, font=\footnotesize\bfseries] {เวกเตอร์การกระจัดลัพธ์ $\vect{d}$} (13.1,3.0);
    \fill[physRoyalBlue] (8.4,1.3) circle (3pt);
    \node[font=\footnotesize, align=center, color=charcoal] at (11.2,0.55) {ระยะทางรวม $s$ ยาวกว่าการกระจัด $\vect{d}$\\แต่มีทิศทางมุ่งหน้าชัดเจน};
\end{tikzpicture}%
}
\caption{การเปรียบเทียบระหว่าง (ก) การเคลื่อนที่แบบบราวน์ที่มีการกระจัดลัพธ์ใกล้ศูนย์ กับ (ข) การว่ายเคลื่อนที่จริงของแบคทีเรียเข้าหาสารอาหาร}
\end{figure}

ในทางตรงกันข้าม เซลล์ที่ไม่มีแฟลกเจลลา (เช่น แบคทีเรียกลุ่มทรงกลม \textit{Staphylococcus aureus} หรือเซลล์ที่ตายแล้ว) เมื่ออยู่ในหยดน้ำภายใต้กล้องจุลทรรศน์จะยังคงสั่นสะเทือนอยู่ตลอดเวลา ปรากฏการณ์นี้เรียกว่า \textbf{การเคลื่อนที่แบบบราวน์} (\textbf{Brownian Motion}) ซึ่งเกิดจากโมเลกุลของน้ำที่เคลื่อนที่ด้วยพลังงานความร้อนวิ่งเข้าชนเซลล์จากทุกทิศทุกทางอย่างสุ่ม ทำให้ผลรวมของเวกเตอร์การกระจัดเมื่อเวลาผ่านไปมีค่าใกล้ศูนย์ ($\vect{d} \approx 0$) เซลล์จึงเพียงแต่สั่นอยู่กับที่โดยไม่เคลื่อนย้ายตำแหน่งไปข้างหน้า

สำหรับการแพร่ของโมเลกุลสารอาหารหรือยาปฏิชีวนะในวุ้นเพาะเชื้อ (เช่น การทดสอบความไวต่อยาปฏิชีวนะด้วยวิธี Disc Diffusion) ระยะทางเฉลี่ยที่โมเลกุลแพร่ไปได้ ($x_{\text{rms}}$) สัมพันธ์กับสัมประสิทธิ์การแพร่ ($D$) และเวลา ($t$) ตามสมการที่ (2.2)

\begin{equation}
    x_{\text{rms}} = \sqrt{2 D t}
\end{equation}

\begin{worked}[การคำนวณอัตราเร็วและความเร็วของแบคทีเรียและการแพร่ของยาปฏิชีวนะ]
\textbf{โจทย์ปัญหา}\quad แบคทีเรีย \textit{Pseudomonas} ตัวหนึ่งว่ายน้ำแบบสลับพุ่งและหมุนสุ่มเป็นระยะทางตามเส้นทางรวม $s = 150\text{ }\mu\text{m}$ ในเวลา $t = 5\text{ s}$ โดยเมื่อวัดระยะห่างแนวตรงจากจุดเริ่มต้นถึงจุดสุดท้ายได้ค่าการกระจัด $d = 90\text{ }\mu\text{m}$ จงหาอัตราเร็วเฉลี่ยและขนาดความเร็วเฉลี่ยของแบคทีเรียตัวนี้

\textbf{วิธีทำทีละขั้นตอนฉบับเข้าใจง่าย}
\begin{enumerate}
    \item \textbf{ขั้นที่ 1 หาอัตราเร็วเฉลี่ยจากระยะทางรวม ($v = s / t$)}
    \[
        v = \frac{150\text{ }\mu\text{m}}{5\text{ s}} = 30\text{ }\mu\text{m/s}
    \]
    \item \textbf{ขั้นที่ 2 หาขนาดความเร็วเฉลี่ยจากขนาดการกระจัด ($|\vect{v}_{\text{avg}}| = d / t$)}
    \[
        |\vect{v}_{\text{avg}}| = \frac{90\text{ }\mu\text{m}}{5\text{ s}} = 18\text{ }\mu\text{m/s}
    \]
\end{enumerate}
\textbf{คำตอบ}\quad แบคทีเรียมีอัตราเร็วเฉลี่ยตามเส้นทางเท่ากับ $30\text{ }\mu\text{m/s}$ และมีขนาดความเร็วเฉลี่ยในแนวตรงเท่ากับ $18\text{ }\mu\text{m/s}$
\end{worked}

\section{ปฏิบัติการประจำสัปดาห์ที่ 2}

\begin{labbox}[ปฏิบัติการที่ 2\quad การเคลื่อนที่ ความเร็ว และการสังเกตการแพร่และ Brownian Motion (3 ชั่วโมง)]
\noindent\textbf{ผลลัพธ์การเรียนรู้ที่ประเมิน}\quad CLO1 (อธิบายเวกเตอร์และอัตราเร็ว), CLO3 (ทักษะการทดลองและบันทึกผลตามจริง), CLO4 (การเขียนกราฟและรายงานกลุ่ม)

\vspace{4pt}
\noindent\textbf{1. วัสดุและอุปกรณ์}
\begin{itemize}
    \item รางทดลองการเคลื่อนที่แนวตรง (หรือหลอดแก้วใสบรรจุของเหลวพร้อมฟองอากาศ/เม็ดบีดส์จำลองเซลล์) และนาฬิกาจับเวลาดิจิทัล
    \item จานเพาะเชื้อใส่น้ำเย็น ($10\,^\circ\text{C}$) และน้ำอุ่น ($45\,^\circ\text{C}$) พร้อมสารละลายสีย้อมเมทิลีนบลู (Methylene Blue) และไม้บรรทัดละเอียด
\end{itemize}

\noindent\textbf{2. ขั้นตอนการทดลอง}
\begin{enumerate}
    \item \textbf{ตอนที่ 1 การวิเคราะห์กราฟระยะทางกับเวลา}\quad ปล่อยฟองอากาศหรือเม็ดบีดส์ให้เคลื่อนที่ในหลอดของเหลว จับเวลาทุกๆ ระยะทาง $10\text{ cm}, 20\text{ cm}, 30\text{ cm}$ และ $40\text{ cm}$ คำนวณอัตราเร็วเฉลี่ย $v = s/t$ และเขียนกราฟความสัมพันธ์ระหว่างระยะทาง (แกนตั้ง) กับเวลา (แกนนอน)
    \item \textbf{ตอนที่ 2 อิทธิพลของอุณหภูมิต่ออัตราเร็วการแพร่}\quad หยดสีย้อมเมทิลีนบลูปริมาตร $20\text{ }\mu\text{L}$ ลงตรงกลางจานเพาะเชื้อที่บรรจุน้ำเย็นและน้ำอุ่น วัดรัศมีวงสีที่แพร่กระจายออกเมื่อเวลาผ่านไป $1, 2$ และ $3$ นาที เปรียบเทียบอัตราเร็วการแพร่
\end{enumerate}

\noindent\textbf{3. ตารางบันทึกผลการทดลอง}
\begin{center}
\footnotesize\setlength{\tabcolsep}{3.5pt}
\begin{tabularx}{\linewidth}{>{\raggedright\arraybackslash}X c c c c}
\toprule
\textbf{ระยะทาง $s$ ($\text{cm}$)} & \textbf{เวลาครั้งที่ 1 ($\text{s}$)} & \textbf{เวลาครั้งที่ 2 ($\text{s}$)} & \textbf{เวลาเฉลี่ย $t$ ($\text{s}$)} & \textbf{อัตราเร็ว $v = s/t$ ($\text{cm/s}$)} \\
\midrule
$10.0\text{ cm}$ & ............ & ............ & ............ & ............ \\
$20.0\text{ cm}$ & ............ & ............ & ............ & ............ \\
$30.0\text{ cm}$ & ............ & ............ & ............ & ............ \\
$40.0\text{ cm}$ & ............ & ............ & ............ & ............ \\
\bottomrule
\end{tabularx}
\end{center}
\end{labbox}

\section*{คำถามท้ายบทที่ 2 (ทบทวนแนวคิดและการคำนวณเบื้องต้น)}
\begin{enumerate}
    \item เพราะเหตุใดการสังเกตเห็นเซลล์แบคทีเรียสั่นสะเทือนอยู่กับที่ภายใต้กล้องจุลทรรศน์จึงยังไม่สามารถสรุปได้ว่าแบคทีเรียชนิดนั้นมีแฟลกเจลลาและเคลื่อนที่ได้จริง จงอธิบายด้วยแนวคิดเวกเตอร์การกระจัด
    \item หากแบคทีเรีย \textit{Vibrio cholerae} สามารถว่ายน้ำในแนวตรงด้วยอัตราเร็วคงที่ $60\text{ }\mu\text{m/s}$ จะต้องใช้เวลากี่วินาทีในการว่ายผ่านเส้นผ่านศูนย์กลางของขอบเขตภาพกล้องจุลทรรศน์ที่กว้าง $300\text{ }\mu\text{m}$
\end{enumerate}
